\documentclass[11pt]{article}
\usepackage{amsmath,amssymb}
\usepackage[margin=1in]{geometry}
\usepackage{booktabs}

\newcommand{\Msun}{M_\odot}
\newcommand{\kpc}{\,\mathrm{kpc}}
\newcommand{\kms}{\,\mathrm{km\,s^{-1}}}
\newcommand{\Gyr}{\,\mathrm{Gyr}}

\title{Tidal disruption of the LMS-1/Wukong progenitor with two embedded
globular clusters and dynamical friction: methods and orbit--anchor diagnostics}
\author{}
\date{}

\begin{document}
\maketitle

\section{Aim}
Yuan et al.~(2020) and Malhan et al.~(2021) associate the globular clusters
NGC\,5024 (M53) and NGC\,5053 with the LMS-1/Wukong accretion event, but their
orbital models contain \emph{no dynamical friction}. We rewind a progenitor
dwarf in the Milky Way (MW) potential, plant the two clusters inside it, and
integrate forward to the present with and without Chandrasekhar friction. We
then confront the resulting stellar debris with the observed LMS-1 stream from
the \textsc{StreamFinder} catalogue (Ibata et al.~2024), and we test how the
result depends on the progenitor halo mass, the disruption time, and crucially
the choice of orbit on which the debris is launched.

\section{Galactic frame and host potential}
All dynamics are integrated in the Milky Way potential of Bovy~(2015,
\texttt{MWPotential2014}), $\Phi_{\rm MW}(\mathbf{x})$. Observed equatorial
phase space is mapped to the Galactocentric frame with
$R_0=8\kpc$, $z_\odot=25\,\mathrm{pc}$, and solar velocity
$\mathbf{v}_\odot=(11.1,\,232.24,\,7.25)\kms$, and back, using the standard
ICRS$\leftrightarrow$Galactocentric transformation. We denote a 6D state
$\mathbf{w}=(\mathbf{x},\mathbf{v})$.

\section{LMS-1 progenitor dark halo}
The progenitor is an NFW halo whose structural parameters follow the
Dutton \& Macci\`o~(2014) concentration--mass relation at formation redshift
$z_{\rm h}=1$. For a virial mass $M_{200}$,
\begin{align}
\log_{10} c_{200} &= a + b\,\log_{10}\!\left(\frac{M_{200}\,h}{10^{12}\,\Msun}\right),\\
a &= 0.520 + (0.905-0.520)\,e^{-0.617\,z^{1.21}}, \qquad b = -0.101 + 0.026\,z,\\
\rho_{\rm crit}(z) &= \rho_{\rm crit,0}\,E^2(z), \qquad
E^2(z)=\Omega_m(1+z)^3+\Omega_\Lambda,\\
r_{200} &= \left(\frac{3 M_{200}}{4\pi\,200\,\rho_{\rm crit}(z)}\right)^{1/3},
\qquad r_s = \frac{r_{200}}{c_{200}},\\
M_{\rm c} &= \frac{M_{200}}{\ln(1+c_{200}) - c_{200}/(1+c_{200})},
\end{align}
with $h=0.70$, $\Omega_m=0.30$, $\Omega_\Lambda=0.70$. The NFW density and
potential are
\begin{equation}
\rho_{\rm NFW}(r)=\frac{M_{\rm c}}{4\pi r_s^3}\frac{1}{(r/r_s)(1+r/r_s)^2},
\qquad
\Phi_{\rm NFW}(r)=-\frac{G M_{\rm c}}{r}\ln\!\left(1+\frac{r}{r_s}\right).
\end{equation}
The stellar body and each cluster are softened Plummer spheres,
\begin{equation}
\Phi_{\rm P}(r)=-\frac{G M}{\sqrt{r^2+b^2}},
\end{equation}
with scale $b_{\rm LMS1}=2\kpc$ for the stars and $b_{\rm GC}=10\,\mathrm{pc}$
for the clusters. The internal circular speed and 1D velocity dispersion of the
dwarf are
\begin{equation}
v_{\rm circ}(r)=\sqrt{r\left|\partial_r\Phi_{\rm LMS1}\right|},
\qquad
\sigma(r)=\max\!\left(\frac{v_{\rm circ}(r)}{\sqrt2},\,\sigma_{\rm floor}\right),
\end{equation}
where $\Phi_{\rm LMS1}=\Phi_{\rm NFW}+\Phi_{\rm P,\,stars}$.

\section{Chandrasekhar dynamical friction}
A body of mass $M_{\rm sat}$ moving with velocity $\mathbf{v}_{\rm rel}$
relative to a background of density $\rho$ and dispersion $\sigma$ feels the drag
\begin{equation}
\mathbf{a}_{\rm DF}
= -\,\frac{4\pi G^2 M_{\rm sat}\,\ln\Lambda\,\rho}{v_{\rm rel}^3}
\left[\operatorname{erf}(X)-\frac{2X}{\sqrt\pi}\,e^{-X^2}\right]\mathbf{v}_{\rm rel},
\qquad
X=\frac{v_{\rm rel}}{\sqrt2\,\sigma}.
\label{eq:chandra}
\end{equation}
Equation~\eqref{eq:chandra} is applied in two places:
\begin{itemize}
\item \emph{halo--MW friction} on the dwarf's orbit, with $\rho=\rho_{\rm MW}$,
$\sigma=\sqrt{\max(-\Phi_{\rm MW}/2,\,\sigma_{\min}^2)}$, $M_{\rm sat}=M_{200}$,
$\ln\Lambda=4.6$;
\item \emph{cluster--halo friction}, with $\rho=\rho_{\rm NFW}$,
$\sigma=\sigma(r)$, $M_{\rm sat}=M_{\rm GC}$,
$\ln\Lambda=\max[2,\ln(M_{200}/M_{\rm GC})]$.
\end{itemize}
The ``no-DF'' runs set $\mathbf{a}_{\rm DF}\equiv 0$ everywhere, reproducing the
Yuan/Malhan assumption.

\section{Time integration}
All trajectories use a kick--drift--kick (KDK) leapfrog with fixed step
$\tau=2^{-11}\Gyr\approx0.49\,\mathrm{Myr}$:
\begin{align}
\mathbf{v}_{n+1/2}&=\mathbf{v}_n+\tfrac{\tau}{2}\,\mathbf{a}(\mathbf{x}_n,\mathbf{v}_n),\nonumber\\
\mathbf{x}_{n+1}&=\mathbf{x}_n+\tau\,\mathbf{v}_{n+1/2},\\
\mathbf{v}_{n+1}&=\mathbf{v}_{n+1/2}+\tfrac{\tau}{2}\,\mathbf{a}(\mathbf{x}_{n+1},\mathbf{v}_{n+1/2}).\nonumber
\end{align}
Backward integration (``rewind'') uses $\tau\to-\tau$.

\subsection{Equations of motion}
The moving halo centre, the (massless) stream tracers, and the live clusters obey
\begin{align}
\mathbf{a}_{\rm centre} &= -\nabla\Phi_{\rm MW}(\mathbf{x}_c)
   + \mathbf{a}_{\rm DF}^{\rm halo},\\
\mathbf{a}_{\rm star} &= -\nabla\Phi_{\rm MW}(\mathbf{x})
   - \nabla\Phi_{\rm NFW}(\mathbf{x}-\mathbf{x}_c)
   - \sum_{i}\nabla\Phi_{{\rm GC},i}(\mathbf{x}-\mathbf{x}_{{\rm GC},i}),\\
\mathbf{a}_{{\rm GC},i} &= -\nabla\Phi_{\rm MW}(\mathbf{x}_i)
   - \nabla\Phi_{\rm NFW}(\mathbf{x}_i-\mathbf{x}_c)
   - \sum_{j\neq i}\nabla\Phi_{{\rm GC},j}(\mathbf{x}_i-\mathbf{x}_j)
   + \mathbf{a}_{{\rm DF},i}^{\rm halo},
\end{align}
where $\mathbf{x}_c(t)$ is the rewound halo trajectory (below) and the stream
tracers are massless.

\section{Progenitor orbit by rewinding the clusters}
The present-day LMS-1 centroid is the \emph{mass-weighted} mean of the two
clusters' measured phase space,
\begin{equation}
\mathbf{w}_{\rm LMS1}(0)=\frac{M_{5024}\,\mathbf{w}_{5024}+M_{5053}\,\mathbf{w}_{5053}}
{M_{5024}+M_{5053}}.
\end{equation}
Mass-weighting is required: a plain average cancels the clusters' opposite
$v_{\rm los}$/$v_z$ and yields an unphysically circular orbit. This state is
rewound for $T$ in $\Phi_{\rm MW}$ ($+\,\mathbf{a}_{\rm DF}^{\rm halo}$ for the DF
run) to give $\mathbf{x}_c(t)$, $t\in[-T,0]$, the moving NFW centre.

\section{Cluster initial conditions (corrected)}
\label{sec:gcic}
Rather than planting the clusters at arbitrary internal offsets, we set their
$t=-T$ states by rewinding their \emph{observed} present-day phase space
$\mathbf{w}_{{\rm GC},i}(0)$ through the full force field (MW $+$ moving NFW $+$
mutual gravity $[+\,$DF$]$). Forward integration then returns each cluster to its
observed position: exactly for the reversible no-DF case, and to
$\lesssim0.1^\circ$ for the (velocity-dependent, weakly irreversible) DF case.
A by-product is that, on its measured orbit, NGC\,5053 lies $\sim89\kpc$
from the LMS-1 centre at $t=-5\Gyr$ (versus $\sim3\kpc$ for NGC\,5024), i.e.\ it
is not tightly co-orbiting the core.

\section{Stellar stream sampling}
$N=10^4$ massless tracers are drawn from a quasi-spherical distribution function
$f(\mathbf{x},\mathbf{v})$ self-consistent with the Plummer stellar density
inside $\Phi_{\rm LMS1}$, then rigidly translated to the rewound centre
$\mathbf{x}_c(-T)$ and integrated forward, feeling MW $+$ moving NFW $+$ both
clusters.

\section{Tidal (Jacobi) radius}
The instantaneous tidal radius of the dwarf at galactocentric distance $R$ is
solved self-consistently from the enclosed masses,
\begin{equation}
r_{\rm t}=R\left[\frac{M_{\rm NFW}(<r_{\rm t})}{3\,M_{\rm MW}(<R)}\right]^{1/3},
\qquad
M_{\rm NFW}(<r)=M_{\rm c}\!\left[\ln\!\Big(1+\tfrac{r}{r_s}\Big)-\frac{r/r_s}{1+r/r_s}\right],
\end{equation}
with $M_{\rm MW}(<R)=v_c^2(R)\,R/G=|a_R|\,R^2/G$.

\section{Observable comparison}
Final 6D debris is transformed to
$(\alpha,\delta,d,\mu_{\alpha\ast},\mu_\delta,v_{\rm los})$. We select catalogue
LMS-1 members, build the convex hull of their $(\alpha,\delta)$ (dilated $15\%$)
as a footprint $\mathcal{F}$, and compare only model particles inside
$\mathcal{F}$. Coverage is the fraction of observed stars with $\geq5$ model
neighbours within $3^\circ$; kinematic agreement uses the median offsets and the
local RMS residuals in $\mu$ and $v_{\rm los}$.

To test whether the debris is launched on the right orbit, we measure the
great-circle distance of each observed star to the projected progenitor orbit
$\{\alpha(t),\delta(t)\}$,
\begin{equation}
\theta_k=\min_t \arccos\big(\hat{\mathbf{n}}_k\cdot\hat{\mathbf{n}}_{\rm orb}(t)\big),
\end{equation}
together with the angle between the on-sky stream direction and the local orbit
tangent.

\section{Stream-fit (Malhan) orbit}
For the literature comparison we also build an IC that matches the
Malhan et al.~(2021) orbit ($r_{\rm peri}\!\approx\!10$,
$r_{\rm apo}\!\approx\!16\kpc$, $e\!\approx\!0.2$,
$L_z\!\approx\!-700\kpc\kms$, inclination $\sim75^\circ$). Anchoring at
apocenter along the observed stream centroid, we solve for the tangential speed
and orientation that reproduce $(r_{\rm peri},r_{\rm apo},L_z)$, verifying by
direct integration.

\section{Experiments and key results}
\begin{itemize}
\item \textbf{DF vs.\ no-DF.} With DF both clusters sink to and are retained in
the core; the rewound centre orbits differ from the no-DF case by up to
$\sim86\kpc$ at $-5\Gyr$ (but negligibly over the last $\sim$Gyr, since both
share the identical present-day state).
\item \textbf{Disruption time $T=5\text{--}9\Gyr$.} The footprint kinematics are
essentially unchanged and the bound fraction drops only mildly with $T$: time is
not the controlling parameter.
\item \textbf{Halo mass $M_{200}=10^{8}\text{--}10^{11}\,\Msun$.} This is the
decisive lever. Light haloes ($\sim10^{8}$) fully disrupt into a long, cold
stream reaching the observed $\alpha\!\approx\!270^\circ$ with the best
proper-motion match but zero bound core; heavy haloes ($\gtrsim10^{10}$) retain
the cluster pair but give a compact, under-extended blob. No single mass both
makes the stream and retains the clusters.
\item \textbf{Orbit anchor.} The cluster-derived orbit lies a median $16^\circ$
off the observed stream (on-sky directions differ by $53^\circ$); debris launched
on it never develops the observed $\mu_\delta$ gradient. The Malhan stream-fit
orbit follows the data to a median $0.7^\circ$; disrupting the same dwarf on it
yields coverage $66\%$, $\alpha_{\max}\!\approx\!270^\circ$, and footprint
$\mu_\delta=-2.37$ vs.\ observed $-2.34\,\mathrm{mas\,yr^{-1}}$. The mismatch was
the launch orbit, not dynamical friction.
\end{itemize}

\end{document}
